\documentclass[10pt,a4paper]{article}
\usepackage[T1]{fontenc}
\usepackage{lmodern}
\usepackage{amsmath,amssymb,amsthm,mathtools}
\usepackage[margin=24mm]{geometry}
\usepackage{microtype}
\usepackage{booktabs,array,longtable}
\usepackage{xcolor}
\usepackage{hyperref}
\usepackage{fancyhdr}
\hypersetup{colorlinks=true,linkcolor=black,citecolor=black,urlcolor=blue,pdfauthor={Matthew J. Colbrook}}
\pagestyle{fancy}
\fancyhf{}
\fancyhead[L]{\small Resolution / MI-21}
\fancyhead[R]{\small 11 September 2026}
\fancyfoot[C]{\thepage}
\setlength{\headheight}{14pt}
\setlength{\parskip}{4pt}
\setlength{\parindent}{0pt}
\newtheorem{theorem}{Theorem}[section]
\newtheorem{proposition}[theorem]{Proposition}
\newtheorem{lemma}[theorem]{Lemma}
\newtheorem{corollary}[theorem]{Corollary}
\theoremstyle{remark}
\newtheorem{remark}[theorem]{Remark}
\newcommand{\M}{\mathbb M}
\newcommand{\C}{\mathbb C}
\newcommand{\R}{\mathbb R}
\newcommand{\tr}{\operatorname{tr}}
\newcommand{\diag}{\operatorname{diag}}
\newcommand{\spec}{\operatorname{spec}}
\newcommand{\rank}{\operatorname{rank}}
\newcommand{\abs}[1]{\lvert #1\rvert}
\newcommand{\norm}[1]{\lVert #1\rVert}
\newcommand{\opnorm}[1]{\lVert #1\rVert_{\infty}}
\newcommand{\schatten}[2]{\lVert #1\rVert_{#2}}
\newcommand{\symmod}{\mathcal S}
\newcommand{\maxmod}{\mathcal M}
\newcommand{\draftnotice}{\begin{center}\small\textit{Proposed mathematical resolution. Complete proof supplied; not submitted or peer reviewed.}\end{center}}

\usepackage{xurl}
\setlength{\emergencystretch}{3em}
\allowdisplaybreaks[1]

\title{MI-21: Resolution}
\author{Matthew J. Colbrook\\
\small Department of Applied Mathematics and Theoretical Physics\\
\small University of Cambridge, Cambridge, United Kingdom\\
\small\href{mailto:m.colbrook@damtp.cam.ac.uk}{m.colbrook@damtp.cam.ac.uk}}
\date{11 September 2026}
\begin{document}
\maketitle
\textbf{Repository status: Solved.}
The complete argument received an independent Codex-agent PASS review
against the exact catalog target.
The original draft was AI-assisted; the reviewed proof below is unchanged.
This is independent agent verification, not external human peer review or
formal proof certification. \href{https://github.com/ajt60gaibb/OpenProblemsInNLA/blob/main/references/colbrook-matrix-2026-09-11/verification/reviews/MI-21-review.md}{Detailed review and scope}.
\par\medskip
% BEGIN REVIEWED PROOF
\section{MI-21: a rational counterexample to the geometric-mean norm conjecture}
\label{sec:mi21}

For positive definite matrices, write
\[
 A\#_tB=A^{1/2}(A^{-1/2}BA^{-1/2})^tA^{1/2},\qquad A\#B=A\#_{1/2}B.
\]
With $A=\sum_iA_i$ and $B=\sum_iB_i$, the target asserts
\begin{equation}\label{eq:mi21-target}
 \left\|\sum_i(A_i^s\#_tB_i^s)^r\right\|
 \leq
 \left\|\left(A^{(1-t)srp/2}B^{tsrp}A^{(1-t)srp/2}\right)^{1/p}\right\|
\end{equation}
for every unitarily invariant norm, $0\leq t\leq1$, and $p,r,s>0$ with $sr\geq1$ \cite[Conjecture 3]{FH2024}.

\begin{theorem}\label{thm:mi21}
Inequality \eqref{eq:mi21-target} is false already for two summands of order two, the operator norm, $s=t=1/2$, and $r=2$. The same counterexample works for every $p>0$ and has rational positive definite inputs satisfying $A=B=I_2$.
\end{theorem}

\begin{proof}
Define the rational matrices
\[
 C=\diag\left(\frac{12}{37},\frac{21}{29}\right),\qquad
 E=\diag\left(\frac{35}{37},\frac{20}{29}\right),\qquad
 S=\frac1{17}\begin{pmatrix}15&8\\8&-15\end{pmatrix}.
\]
They satisfy $C,E>0$, $C^2+E^2=I_2$, and $S=S^*=S^{-1}$. Set
\[
 A_1=C^2,\quad A_2=E^2,\qquad
 B_1=SE^2S,\quad B_2=SC^2S.
\]
All four inputs are rational and positive definite. Also
$A_1+A_2=B_1+B_2=I_2$. Thus the matrix inside the norm on the right of \eqref{eq:mi21-target} is exactly $I_2$, for every $p>0$.

Put $D=SES$ and $G=C\#D$. Symmetry and unitary covariance of the geometric mean imply
\[
 E\#(SCS)=SGS.
\]
For the chosen parameters, the left-hand matrix is therefore
\[
 L=G^2+SG^2S.
\]
We compute it exactly. For any $2\times2$ positive definite $C,D$, if
\[
 \delta=\sqrt{\det D/\det C},\qquad h=\tr(C^{-1}D)+2\delta,
\]
then
\begin{equation}\label{eq:mi21-2by2}
 C\#D=\frac{D+\delta C}{\sqrt h},\qquad
 (C\#D)^2=\frac{(D+\delta C)^2}{h}.
\end{equation}
Indeed, applying the Cayley--Hamilton identity to $C^{-1/2}DC^{-1/2}$ shows that its positive square root is
\[
 \frac{C^{-1/2}DC^{-1/2}+\delta I}{\sqrt{\tr(C^{-1}D)+2\delta}}.
\]
Congruence by $C^{1/2}$ proves \eqref{eq:mi21-2by2}.

In the present example,
\[
 \delta=\frac53,\qquad
 h=\frac{61697295}{8682716},\qquad
 D+\delta C=\frac1{310097}
 \begin{pmatrix}443355&33000\\33000&605715\end{pmatrix}.
\]
Rational matrix multiplication in \eqref{eq:mi21-2by2} gives
\[
 G^2=\frac1{12158163}
 \begin{pmatrix}3516940&616000\\616000&6547660\end{pmatrix},
\]
and hence
\begin{equation}\label{eq:mi21-L}
 L=\frac1{12158163}
 \begin{pmatrix}8216600&-985600\\-985600&11912600\end{pmatrix}.
\end{equation}
For $w=(1,-4)^T$, direct multiplication yields
\[
 Lw=\frac{1351000}{1350907}w,
 \qquad\frac{1351000}{1350907}=1+\frac{93}{1350907}>1.
\]
Therefore $\opnorm{L}>1=\opnorm{I_2}$, contradicting \eqref{eq:mi21-target}.
\end{proof}

\begin{remark}[Exact certificate and scope]
The counterexample is not based on a rounded numerical eigenvalue. Every displayed matrix and every comparison after \eqref{eq:mi21-2by2} is rational. Its other eigenvalue is $7970200/12158163<1$. The parameter condition is met with $sr=1$. No claim is made here about the restricted parameter regimes already established in the cited paper.
\end{remark}

% END REVIEWED PROOF
\begin{thebibliography}{1}

\bibitem{FH2024}
Shaima'a Freewan and Mostafa Hayajneh.
\newblock On a conjecture related to the geometric mean and norm inequalities.
\newblock {\em Mathematical Inequalities \& Applications}, 27(1):193--200,
  2024.
\newblock Conjecture 3, page 195.
\newblock URL: \url{https://files.ele-math.com/articles/mia-27-15.pdf}, \href
  {https://doi.org/10.7153/mia-2024-27-15} {\path{doi:10.7153/mia-2024-27-15}}.

\end{thebibliography}

\end{document}
